\documentclass[11pt]{article} 

% Core math and formatting packages to replace evan.sty
\usepackage{amsmath}
\usepackage{amssymb}
\usepackage{amsthm}
\usepackage{geometry}
\geometry{margin=1in} % Standard readable margins
\usepackage{fancyhdr} % For custom headers

\usepackage{dsfont}  
\usepackage[most]{tcolorbox}
\usepackage{multicol}
\usepackage{comment}

% Define custom environments that were previously handled by evan.sty
\newtheorem{theorem}{Theorem}[section]
\newtheorem*{theorem*}{Theorem}
\newtheorem{exercise}{Exercise}[section]

\theoremstyle{definition} % Makes the text Roman instead of italics for examples
\newtheorem{example}{Example}[section]

\definecolor{midnightblue}{RGB}{25, 25, 112}
% Define a custom box for your problems
\newtcolorbox{probbox}[1][]{
    enhanced,
    breakable,
    colback=blue!5!white,
    colframe=midnightblue, 
    fonttitle=\bfseries,
    arc=2mm,
    title=Problem Statement
    #1
}

\title{A Perron-Theoretic Proof of the Polynomial Goldbach Theorem}
\author{Ayush Mandal (B.Math 2nd Year) \\ \normalsize Indian Statistical Institute, Bangalore}
\date{}

\begin{document}

\maketitle

% Header setup using fancyhdr instead of KOMA-script
\pagestyle{fancy}
\fancyhf{} % Clears all default headers and footers
\lhead{\small \textbf{Ayush Mandal(B.Math 2nd Year)}} % Left header
\cfoot{\thepage} % Page number at the bottom center
\renewcommand{\headrulewidth}{0.4pt} % Adds a subtle line under the header

\begin{abstract}
    In 1965, David Hayes showed that any polynomial $f\in \mathbb{Z}[X]$ of degree $n$ can be expressed as 
    the sum of two irreducible integer polynomials of degree $n$. In general it is true for any Noetherian ring $R$
    with infinitely many maximal R-ideals which was proved in the $21st$ century. Here I have used Perron's criterion
    of irreducibility to show the result for $R=\mathbb{Z}$.
\end{abstract}

\section{Introduction}
We all have definitely heard of the famous \textit{Goldbach Conjecture}.
It states that any even positive integer greater than $2$ can be expressed as sum of two prime numbers. 
 We say an element $a$ in a ring $R$ irreducible, if 
$a=bc$ implies either $b$ or $c$ is a unit (which has multiplicative inverse). In a Principal Ideal Domain (PID), every irreducible element is prime 
whose proof one can find in any algebra book. Since $\mathbb{Z}$ is a PID, every irreducible element is prime. Thus we can refine our statement of Goldbach 
conjecture by replacing prime with irreducible elements of $\mathbb{Z}$. Getting motivated from this problem, we ask the same question in polynomial ring.
Can every polynomial over any arbitrary ring of degree $n$ be expressed as the sum of irreducible polynomials of the same ring having degree $n$? The answer is not 
always. But for Noetherian rings having infinitely many maximal ideals we can. Since $\mathbb{Z}$ is such a ring, for any $f\in \mathbb{Z}[X]$ there exists 
irreducible polynomials $g,h\in \mathbb{Z}[X]$ such that $f=g+h$ with $\deg(f)=\deg(g)=\deg(h)$. David Hayes proved it using Eisenstein's criterion of 
irreducibility. Here I will be proving this theorem using a very strong criterion for irreducibility called Perron's criterion which I will be proving.

\section{Preliminaries}
For this paper, one need to know some Complex analysis as the well known \textit{Rouché's theorem} will be used in the proof of the Perron's criterion. I will not be proving this theorem
as one can easily find it in any standard complex analysis book.
\newpage

\begin{theorem}[\textit{Rouché's theorem}]
    Let f and g be analytic functions on and inside a simple closed curve C. Suppose
that $|f(z)| > |g(z)|$ for all points $z$ on $\mathcal{C}$. Then $f$ and $f + g$ have the same number of zeros (counting
multiplicities) interior to $\mathcal{C}$.
\end{theorem}

\section{Perron's Irreducibility Criterion}
\begin{theorem}[Perron]
    Let $P(x) = x^n + a_{n-1} x^{n-1} +\cdots+ a_1 x + a_0$ be an integer polynomial with $a_0\neq 0$ and
    $$|a_{n-1}|>1+|a_{n-2}|+\cdots+|a_1|+|a_0|$$
Then $P(x)$ is irreducible.
\end{theorem}
\begin{proof}
    Let $\alpha\in \mathbb{C}$ be a root of $P(x)$ with $|\alpha|=1$. Then
    $$-a_{n-1}\alpha^{n-1}=1+\alpha^n+a_{n-2}\alpha^{n-2}+\cdots+a_0$$
    $$|a_{n-1}|\le 1+ |a_{n-2}|+\cdots +|a_1|+|a_0|$$ which contradicts the condition. Thus, 
    $|\alpha|\neq 1$ or no roots lie on the unit circle. Now consider polynomials $f(x)=a_{n-1}x^{n-1}$ and $g(x)=P(x)-a_{n-1}x^{n-1}$ and the simple closed curve 
    $\mathcal{C}=\{z\in \mathbb{C}: |z|=1\}$. By the given condition we can easily say that 
    $$|f(z)|=|a_{n-1}z^{n-1}|>|P(z)-a_{n-1}z^{n-1}|=|g(z)|$$ for all $z\in \mathcal{C}$. Thus by  Rouché's theorem, we can say that $a_{n-1}z^{n-1}$ and $P(z)$ has same number of zeroes inside $\mathcal{C}$. 
    Since $f$ has exactly $n-1$ zeroes inside, $P(z)$ also has exactly $n-1$ zeroes strictly inside $\mathcal{C}$ (as no roots can lie on $\mathcal{C}$) and another say $\alpha $ with $|\alpha|>1$. 
    Now suppose $P$ is reducible, then there exists $f,g\in \mathbb{Z}[X]$ with $\deg(f),\deg(g)\ge 1$ such that $$P(x)=f(x)g(x)$$ Since $P(0)=a_0\neq 0$, $|f(0)|,|g(0)|\ge 1$.
    WLOG, let $\alpha$ be a root of $g(x)$. Let $\deg(f)=k<n$ and $\alpha_1,\alpha_2,\cdots, \alpha_k$ be the roots of $f(x)$. Then $|\alpha_i|<1$ for all $1\le i\le k$. Also since $P(x)$ is a monic polynomial, $f,g$ also has to be
    monic. Using Vieta's theorem we can say that  $$1\le |f(0)|=\left|\prod_{i=1}^k\alpha_i\right|<1$$ which is a contradiction. Thus, $P$ is irreducible.
\end{proof}
\newpage

This criterion makes our life very easy while solving irreducibility problems. One can solve the following example problem directly by invoking the theorem.
\begin{example}[IMO 1993]
     Let $f(x) = x^n + 5x^{n-1} + 3$, where $n > 1$ is an integer. Prove that $f(x)$ cannot be expressed
    as the product of two nonconstant polynomials with integer coefficients.
\end{example}
The following are some direct exercises to enjoy!
\begin{exercise}
Let $P(z)=a_nz^n+a_{n-1}z^{n-1}+\cdots+a_0$ be a polynomial with $a_i\in \mathbb{C}$ such that 
$$2|a_k|>\sum_{i=0}^n|a_i|$$ for some $0\le k\le n$. Then show that exactly $k$ many roots of $P$ are strictly inside the unit circle and the 
other $n-k$ outside the circle.
\end{exercise}

\begin{exercise}
    Let $P(z)=a_nz^n+a_{n-1}z^{n-1}+\cdots+a_0$ be a polynomial with $a_i\in \mathbb{Z}$ and $|a_0|$ is prime such that 
    $$2|a_0|>\sum_{i=0}^n|a_i|$$ Show that $P$ is irreducible in $\mathbb{Z}[X]$.
\end{exercise}

\begin{exercise}[Iran 2003]
     Let $f_1,f_2,\cdots,f_n$ be polynomials with integer coefficients. Show that there exists a reducible
     polynomial $g(x)\in \mathbb{Z}[x]$ such that $f_i(x)+g(x)$ is irreducible for $i=1,2,\cdots,n$.
\end{exercise}

\section{The Main Theorem}
\begin{theorem*}
    For any polynomial $P\in \mathbb{Z}[X]$ of $\deg(P)\ge 1$, there exists $f,g\in \mathbb{Z}[X]$ irreducible polynomials with $\deg(P)=\deg(f)=\deg(g)$ such that 
    $$P(x)=f(x)+g(x)$$
\end{theorem*}
\begin{proof}
    In my proof I will be explicitly constructing polynomials $f,g$ to get the theorem. Suppose 
    $P(x)=a_nx^n+a_{n-1}x^{n-1}+\cdots+a_0$ with $a_i\in \mathbb{Z}[X]$. Define $$f(x)=b_nx^n+(a_{n-1}-pM) x^{n-1}+a_{n-2}x^{n-2}+\cdots+a_1x+(a_0-p)$$ $$g(x)=c_nx^n+pMx^{n-1}+p$$ where 
    $p$ is some prime such that $p\nmid c_n$ and $a_0-p\neq 0$.\newpage If $a_n=1$, $b_n=-1$ and $c_n=2$ else $b_n=1$ and $c_n=a_n-1$. Now choose $M$ such that $p\nmid M$ and 
    $$|a_{n-1}-pM|>1+|a_{n-2}|+\cdots+|a_1|+|a_0-p|$$ Thus one can easily check that $P(x)=f(x)+g(x)$. Observe that $g$ is an Eisenstein polynomial and $f$ satisfies the Perron's criterion as per the construction, thus both $f,g$ are irreducible.
    This completes the proof!
\end{proof}

\begin{example}
    Consider a polynomial $p(x)=x^5+4x^4-2x^2+7$. Using the above construction, we take the polynomials $g(x)=2x^5+3\cdot 8x^4+3$ and $f(x)=-x^5-20x^4-2x^2+4$ where both $f,g$ are
    irreducible and $p=f+g$.
\end{example}

The one who are interested to explore more on this, can use the following references to get the proof of the general theorem.

\begin{thebibliography}{9} % '9' means you have fewer than 10 references. Use '99' for up to 99.

\bibitem{hayes1965}
D. R. Hayes, \textit{A Goldbach Theorem for Polynomials with Integral Coefficients}, 
The American Mathematical Monthly, vol. 72, no. 1, 1965, pp. 45--46.

\bibitem{pollack2011}
P. Pollack, \textit{On Polynomial Rings with a Goldbach Property}, 
The American Mathematical Monthly, vol. 118, no. 1, 2011, pp. 71--77.

\end{thebibliography}
\end{document}